\scanpage{210}
\mspage{2}
% archive sheet number: 209

(ii) On a les conditions G0 à G5

G0 $C$ est une $\mathcal U$-catégorie, l'ens des \deleted{\unsure{classes mod}} isomorphies d'objets de $C$ est $\mathcal U$-petit

G1 Existence de $\underleftarrow{\mathrm{Lim}}$ \inserted{finies} et \deleted{$\underrightarrow{\mathrm{Lim}}$}\inserted{sommes} \uline{finies}

G2 \deleted{Les \unsure{équi|épim.}relations d'équivalence sont effectives,} \deleted{les épimorphismes sont effectifs universels,} les sommes finies \deleted{sont} disjointes \deleted{\unsure{eff|effect|épi.}}\unsure{universelles|effectives}.

G3 Tout morphisme $X\longrightarrow Y$ se factorise en \deleted{épimorphisme}\inserted{morphisme de descente effective} suivi de isom. sur sommand direct.

G4 \deleted{Tout} L'objet \inserted{final} est \deleted{\unsure{localt|une} (ou au choix, \unsure{arbitraire})} somme \deleted{\unsure{disjoint|disjointe}} effective univ. d'objets connexes.

G5 \deleted{pour les \unsure{sous-objets|sous - objets} [ou encore : l'objet final est} \deleted{somme \uline{finie} d'objets connexes].} soit $d:X\longrightarrow Y$ avec $X,Y$ connexes, alors $\exists$ entier $N$ tel que pour tt $Y'\longrightarrow Y$, $Y'\ne\emptyset_C$, \deleted{\unsure{l'élément}} $X\times_Y Y'$ ne contient pas $Y'\times\unsure{[1,N+1]|[1,N+1[}$ \inserted{comme sommand}.

(\deleted{\unsure{vi}}\inserted{iii}) \deleted{Il existe} $C$ est équivalente à une catégorie produit \deleted{\illegible{1}} $\prod_{i\in I}C_i$, les $C_i$ \uline{galoisiennes}.

[NB les $C_i$ correspondent donc aux composantes connexes \inserted{$e_i$} de l'objet final ; la donnée d'une équivalence $C_i\simeq C(\pi_i)$ correspond à la donnée d'un foncteur fibre \inserted{de $C/e_i$}, ou d'un ``rev. universel'' de l'objet \deleted{\unsure{$f$}} $e_i$].

\uline{Démonstration de A}

$(\mathrm{iii})\Longrightarrow(\mathrm{ii\ bis})\Longrightarrow(\mathrm{ii})\Longrightarrow(\mathrm{i})$ triviaux, et $(\mathrm{i})\Longrightarrow(\mathrm{iii})$ aussi, car si $X\in\mathscr T_{\unsure{\mathrm{revf}|\mathrm{rev}\,f|\mathrm{cf}}}$, alors l'objet final $e$ de $\mathscr T$ est somme (dans $\mathscr T$) d'objets $e_i$ $(i\in\mathbb N)$ avec $X|e_i$ de rang $i$. De plus, $\underline{\mathrm{Isom}}_{e_i}([1,i]_{e_i},X)=P_i$ est alors un objet de $\mathscr T_{\mathrm{revf}}$ tel que $\unsure{P|R}_i\longrightarrow e_i$ épim (c'est un torseur sous $(\mathfrak S_i)_{e_i}$ !) et tel que

